%% Proof annex fragment for Part 2.1 calculi.  Complete proofs only; every
%% item has exactly one body back-reference.

\begin{annexitem}{annex:p2.calc.1}{thm:p02-nd-open-substitution}{Substitution for open assumptions}
\begin{proof}
For designated occurrences reserve pairwise disjoint label sets also
disjoint from \(D\)'s labels; consistently rename their copies of \(E\) into
them.  Induct on \(D\)'s height.  An assumption becomes its reserved copy iff
designated.  At an inference, inductively transform each premise and reapply
the rule with its original discharge labels.  Designated occurrences are
open, so none is discharged there; each original discharge remains present
and nonvacuous.  Renaming preserves internal discharges, while freshness
prevents \(D\)'s rules from discharging copy assumptions.  The reconstructed
derivation concludes \(\psi\) and has exactly the stated open assumptions.
\end{proof}
\end{annexitem}

\begin{annexitem}{annex:p2.calc.2}{thm:p02-hilbert-deduction}{Deduction theorem for the Hilbert calculus}
\begin{proof}
First suppose
\(\Gamma\derives_{\mathsf H_{\mathsf{CPL}}}\varphi\to\psi\).
The same derivation is available from \(\Gamma\cup\{\varphi\}\); append
\(\varphi\) as a premise and then apply modus ponens.  Hence
\(\Gamma\cup\{\varphi\}\derives_{\mathsf H_{\mathsf{CPL}}}\psi\).

Conversely, let \(\theta_0,\ldots,\theta_n=\psi\) be a Hilbert derivation
from \(\Gamma\cup\{\varphi\}\).  Induct on \(i\) to derive
\(\varphi\to\theta_i\) from \(\Gamma\).

If \(\theta_i\in\Gamma\), or \(\theta_i\) is an axiom-schema instance,
derive \(\theta_i\), take the A1 instance
\(\theta_i\to(\varphi\to\theta_i)\), and apply modus ponens.  If
\(\theta_i=\varphi\), the following A1 and A2 instances derive
\(\varphi\to\varphi\):
\[
\begin{aligned}
 &\varphi\to((\varphi\to\varphi)\to\varphi),\\
 &\varphi\to(\varphi\to\varphi),\\
 &[\varphi\to((\varphi\to\varphi)\to\varphi)]
   \to([\varphi\to(\varphi\to\varphi)]
   \to(\varphi\to\varphi)).
\end{aligned}
\]
Apply modus ponens first to the first and third displayed lines, and then to
the result and the second line.

Finally suppose \(\theta_i\) follows by modus ponens from earlier
\(\eta\) and \(\eta\to\theta_i\).  The induction hypothesis gives
\(\varphi\to\eta\) and \(\varphi\to(\eta\to\theta_i)\).  Use the A2
instance
\[
 [\varphi\to(\eta\to\theta_i)]
 \to([\varphi\to\eta]\to[\varphi\to\theta_i])
\]
and modus ponens twice to obtain \(\varphi\to\theta_i\).  These exhaust
the ways a Hilbert line is licensed.  At \(i=n\) the result is
\(\Gamma\derives_{\mathsf H_{\mathsf{CPL}}}\varphi\to\psi\).
\end{proof}
\end{annexitem}
